% =====================================================================
%  Davidson Student Scholars Proposal -- AY2026-27
%  Thuc Bao Tran, SJSU College of Engineering
%
%  Compiles with pdflatex (no unicode characters used in the body; all
%  special glyphs are LaTeX commands, so it also builds under xelatex
%  and lualatex without changes).
%
%    pdflatex PROPOSAL_v3.tex
%
%  Page budget: DSS allows <= 3 body pages plus cover/reference pages.
%  If the body runs long, the levers in order of least damage are:
%    1. \geometry margins (currently 0.72in)
%    2. the 9.6pt/11.6pt size in \documentclass + \linespread
%    3. \small -> \footnotesize on the tables
%  Trim typography before trimming content.
% =====================================================================
\documentclass[10pt,letterpaper]{article}

\usepackage[T1]{fontenc}
\usepackage[utf8]{inputenc}
\usepackage{lmodern}
\usepackage[margin=0.72in,top=0.62in,bottom=0.72in]{geometry}
\usepackage{booktabs}
\usepackage{array}
\usepackage{enumitem}
\usepackage{titlesec}
\usepackage{xcolor}
\usepackage{mdframed}
\usepackage{microtype}
\usepackage[hidelinks]{hyperref}

\definecolor{dss}{HTML}{0D2B33}
\definecolor{dssrule}{HTML}{0D6672}
\definecolor{dssbg}{HTML}{F4F8F8}

% --- compact sectioning; headings are lettered manually to match the
%     DSS A-E structure rather than arabic \section numbering ----------
\titleformat{\section}{\normalfont\large\bfseries\color{dss}}{}{0pt}{}
\titlespacing*{\section}{0pt}{9pt plus 2pt minus 2pt}{4pt}
\titleformat{\subsection}{\normalfont\normalsize\bfseries\color{dss}}{}{0pt}{}
\titlespacing*{\subsection}{0pt}{7pt plus 1pt minus 1pt}{3pt}

\setlist[itemize]{leftmargin=1.1em,itemsep=1pt,topsep=2pt,parsep=0pt}
\setlist[enumerate]{leftmargin=1.4em,itemsep=1.5pt,topsep=2pt,parsep=0pt}

\linespread{1.02}
\setlength{\parskip}{3.2pt}
\setlength{\parindent}{0pt}
\setlength{\tabcolsep}{5pt}
\renewcommand{\arraystretch}{1.06}

% shorthands for units that recur throughout
\newcommand{\apthree}{AP\textsubscript{3D}}
\newcommand{\apbev}{AP\textsubscript{BEV}}
\newcommand{\dg}{\ensuremath{^\circ}}
\newcommand{\too}{\ensuremath{\rightarrow}}

\pagestyle{plain}

\begin{document}

% ------------------------------------------------------------------ head
{\footnotesize
\textbf{Name:} Tran, Thuc Bao \quad$\cdot$\quad
\textbf{Email:} \href{mailto:thucbao.tran@sjsu.edu}{thucbao.tran@sjsu.edu}\\
\textbf{Department:} Computer Engineering, College of Engineering \quad$\cdot$\quad
\textbf{Faculty Mentor:} Prof.\ Kaikai Liu\\
\textbf{Track:} Undergraduate Research Project (\$1,500) \quad$\cdot$\quad
\textbf{Period:} Fall 2026 -- Spring 2027 \quad$\cdot$\quad
\textbf{Final report due:} 30 May 2027\par}

\vspace{5pt}
{\color{dss}\large\bfseries Which Component Fails When a Roadside Detector Moves?\par}
\vspace{3pt}

% -------------------------------------------------------------- synopsis
\begin{mdframed}[linewidth=0pt,leftline=true,linecolor=dssrule,
                 innerleftmargin=9pt,innerrightmargin=7pt,
                 innertopmargin=6pt,innerbottommargin=6pt,
                 backgroundcolor=dssbg,skipabove=5pt,skipbelow=7pt]
{\bfseries\color{dss}Synopsis of Proposed Research} \hfill {\footnotesize (240 words)}

\smallskip
Roadside 3D detectors are trained at one intersection and deployed at another. Published
accuracy is measured on the training distribution; deployment is not. The loss is known to be
large --- a reported 70--90\% drop in detection rate across changes in lidar, geography and
weather \cite{ms3d} --- but the field reports it as a single number and attempts to close it
wholesale.

Preliminary measurements on a Jetson Orin Nano suggest the loss is not uniform. Evaluating
NVIDIA's pretrained V2XFusion \cite{nvidia} zero-shot from DAIR-V2X-I (Beijing) \cite{dair} to
TUMTraf Intersection (Munich) \cite{tumtraf}, detection and position largely survive while
heading estimation collapses: vehicles within $\pm 5\dg$ of true heading fall from 88\% to
28\%, with the median error unchanged at $-0.2\dg$ in both domains. Because roughly $20\dg$ of
heading error on a $4.1 \times 1.9$\,m box clears IoU 0.25 but not IoU 0.5, this one mode
explains why Car accuracy falls 69.70 \too\ 0.13 AP while Pedestrian and Cyclist retain 56\%
and 65\%.

Heading is already treated as a separate error axis in-domain --- nuScenes reports mAOE
independently of mATE and mASE \cite{nuscenes} --- but a systematic search of the adaptation
literature found it isolated only in-domain or under adversarial perturbation, never under
natural cross-dataset shift, and never for roadside sensors.

This project will test whether domain fragility concentrates in specific components, and
whether the fragile one can be repaired without target labels. Four label-free interventions will
be measured against a supervised oracle.
\end{mdframed}

% =================================================================== A
\section{A.\quad Introduction}

\subsection{A.1\quad Current Research}

\textbf{The cross-dataset drop is large and label-free adaptation partly closes it.}
Self-training recovers 16--75\% of the source-only-to-oracle gap across four LiDAR transfers
using no target labels and no target statistics; on Waymo\too KITTI it lifts Car \apthree\ from
27.48 to 61.83 against a 73.45 oracle \cite{st3d}. Cheaper still, a purely source-side
augmentation --- random object scaling --- recovers about 86\% of what the weakly supervised
size prior buys ($+27.19$ \apthree\ against $+31.72$), and \emph{outperforms} it on cyclist
\cite{st3dpp}.

\textbf{Weak supervision is not automatically better.} Statistical Normalization consumes
target-domain object-size statistics yet degrades transfer when the size gap is small:
nuScenes\too KITTI \apbev\ falls 51.84 \too\ 40.03, a $-37.55\%$ closed gap \cite{st3d}.

\textbf{Orientation is already an established, separately-reported error axis.} nuScenes
defines mAOE as a standalone true-positive error alongside mATE and mASE, and computes mASE
only after aligning orientation, decoupling scale from heading by construction
\cite{nuscenes}. Earlier work defines paired full- and half-range metrics,
$\mathrm{FOE} = |(\theta - \hat{\theta}) \bmod 360\dg|$ and
$\mathrm{HOE} = |(\theta - \hat{\theta}) \bmod 180\dg|$, so that $180\dg$ flips are
distinguishable from angular imprecision \cite{cui}.

\textbf{Why heading is hard has a published explanation.} Cui et al.\ observe that ``the front
and back of a vehicle may not be easily distinguishable from the LiDAR point cloud,'' and that
parked vehicles fail because they ``have no moving trajectory predictions that could be used to
reliably infer the orientations'' \cite{cui}. Heading is under-determined by appearance;
detectors lean on motion to resolve it.

\subsection{A.2\quad Limitations of Current Research}

\textbf{Every adaptation result above is vehicle-mounted and LiDAR-only.} Waymo, KITTI,
nuScenes and Lyft are all ego-vehicle datasets. A roadside sensor is static, has no ego-motion,
and in the deployed single-frame configuration has no trajectory cue at all --- precisely the
cue detectors are shown to depend on for heading \cite{cui}.

\textbf{Heading has never been isolated under natural domain shift.} It is measured in-domain
(mAOE, FOE/HOE) and under adversarial perturbation, where error decomposition shows yaw is
disproportionately sensitive and mAP-style metrics hide it \cite{robust}. MS3D++ treats heading
as a class-dependent pseudo-label problem --- catastrophic for elongated vehicles because it
destroys IoU, tolerable for BEV-symmetric pedestrians --- but reports no yaw-specific metric
\cite{ms3d}. None of this is cross-dataset.

\textbf{Adaptation is evaluated as one number.} No work reports which \emph{component} failed,
so it is unknown whether the loss is diffuse or concentrated, and therefore unknown whether a
targeted, cheap repair is even possible.

\textbf{A methodological trap is documented.} Self-training's headline result is
model-selection sensitive: Easy 3D AP has been reported fluctuating between 27.9\% and 60.9\%
across randomizations, because best-epoch selection leaks target information
\cite{st3drepro}. A method that is label-free during training can become label-dependent at
checkpoint selection.

\subsection{A.3\quad Research Gap}

A systematic, adversarially-verified literature search returned \textbf{no} cross-dataset
results for roadside datasets (DAIR-V2X-I, Rope3D, TUMTraf), \textbf{no} LiDAR+camera fusion
transfers, and \textbf{nothing} on INT8 quantization versus out-of-distribution robustness.
That is a limit of the search rather than proof of absence --- the roadside literature exists
--- but it establishes that per-component fragility for roadside 3D detection is, at minimum,
not a well-covered question.

% =================================================================== B
\section{B.\quad Specific Proposed Research and Why Important}

\textbf{Question.} Does domain fragility in roadside 3D perception concentrate in specific
network components, and can the fragile component be repaired using only unlabeled target data?

\textbf{Preliminary evidence, already collected.} A full pipeline was built and validated on a
borrowed Jetson Orin Nano. It reproduces NVIDIA's published DAIR-V2X-I accuracy
\cite{dair,nvidia} to within 0.02--0.04 AP on Car and Pedestrian, establishing that subsequent
measurements are trustworthy rather than artifacts. Zero-shot transfer to TUMTraf Intersection
\cite{tumtraf} (2,160 frames) gives:

\begin{center}\small
\begin{tabular}{@{}lccc@{}}
\toprule
\textbf{3D AP, moderate} & \textbf{DAIR-V2X-I} & \textbf{TUMTraf} & \textbf{retained} \\
\midrule
Car @ IoU 0.5          & 69.70 & 0.13  & 0.2\%  \\
Pedestrian @ IoU 0.25  & 49.39 & 27.61 & 55.9\% \\
Cyclist @ IoU 0.25     & 57.90 & 37.76 & 65.2\% \\
\bottomrule
\end{tabular}
\hfill
\begin{tabular}{@{}lcc@{}}
\toprule
\textbf{Heading error} & \textbf{DAIR-V2X-I} & \textbf{TUMTraf} \\
\midrule
median (signed)    & $-0.2\dg$ & $-0.2\dg$ \\
standard deviation & $12.4\dg$ & $24.0\dg$ \\
within $\pm 5\dg$  & 88\%      & 28\%      \\
\bottomrule
\end{tabular}
\end{center}

The median is identical across domains, which excludes a coordinate or calibration error: the
bias is correct and the reliability is not. The error is multi-modal, clustering near
$+15\ldots30\dg$ and $-30\ldots-60\dg$ --- consistent with a learned prior over road directions
rather than random degradation. The class pattern independently matches MS3D++'s observation
that heading error destroys IoU for elongated vehicles while symmetric classes tolerate it
\cite{ms3d}.

\textbf{Why it matters.} A missed detection is a failure mode downstream planners are built to
tolerate. A confident \emph{wrong heading} is not: heading feeds motion prediction, so a
correctly located vehicle with a $25\dg$ heading error yields a predicted trajectory going
somewhere the vehicle is not. If fragility is concentrated, repair can be local and label-free;
if diffuse, retraining on labeled target data is the only option, which does not scale to
per-intersection deployment.

% =================================================================== C
\section{C.\quad Methodology and Why Innovative}

\textbf{C.1 Per-component fragility profile.} Adapt the sub-task attribution protocol of
monodle \cite{monodle}: substitute ground truth for one predicted quantity at a time ---
heading, center, size, class --- and re-measure AP. Applied both in-domain and cross-domain,
the \emph{difference} in each substitution's effect attributes the drop to specific components.
monodle applies this in-domain to a monocular detector; applying it across a natural domain
gap, to a LiDAR+camera roadside detector, is what is new.

\textbf{C.2 Test the road-geometry hypothesis.} If heading rests on a learned road-angle prior,
the error clusters must align with the \emph{source} intersection's lane bearings. This is
directly falsifiable, and is scheduled first because a null result redirects the work.

\textbf{C.3 Four label-free repairs, measured against an oracle.}

\begin{center}\small
\begin{tabular}{@{}p{0.36\textwidth}p{0.22\textwidth}p{0.34\textwidth}@{}}
\toprule
\textbf{Intervention} & \textbf{Supervision} & \textbf{Precedent} \\
\midrule
Recalibrate INT8 ranges on unlabeled target frames & none --- forward passes &
  untested in the verified literature \\
Adapt orientation head only, backbone frozen & pseudo-labels & --- \\
Random object scaling at source pre-training & none --- source-side &
  $+27.19$ \apthree\ on car \cite{st3dpp} \\
Local-structure input features & none --- source-side &
  $>$21 mAP Waymo\too KITTI \cite{gblobs} \\
\bottomrule
\end{tabular}
\end{center}

A supervised oracle establishes the recoverable ceiling so each repair is reported as
\emph{fraction of gap closed}, the convention used by the self-training literature \cite{st3d}.

\textbf{C.4 Guard against the documented traps.} Checkpoints will be selected on a
source-domain validation split, never on target performance, because best-epoch selection on
the target leaks label information and inflates results \cite{st3drepro}. Statistical
Normalization will be included as a \emph{negative control} rather than a baseline to beat,
since it is known to degrade transfer when the size gap is small \cite{st3d}. Per-class AP will
be reported throughout, never a single averaged mAP, because the class balance differs between
the two domains.

\textbf{C.5 Compression as an axis, not the headline.} Each repair is evaluated at FP16 and
INT8. Preliminary data shows quantization adds $0.0\dg$ of heading error in-domain and
$+1.2\dg$ out-of-domain; establishing whether that is real requires repeated calibration seeds
and per-class confidence intervals, which C.3 supplies.

\textbf{Why innovative.} Existing adaptation methods treat the detector as one object and ask
how much AP returns. This asks \emph{which part broke}, then repairs that part. It also moves
the orientation question out of the vehicle-mounted, multi-frame setting where it was
identified and into the static, single-frame roadside setting where the motion cue that
previously rescued it does not exist.

% =================================================================== D
\section{D.\quad Milestones and Timeline}

\begin{center}\small
\begin{tabular}{@{}p{0.14\textwidth}p{0.47\textwidth}p{0.32\textwidth}@{}}
\toprule
\textbf{Period} & \textbf{Milestone} & \textbf{Deliverable} \\
\midrule
Sep--Oct 2026 & Lane-bearing correlation (C.2); second target domain converted &
  Hypothesis confirmed or refuted \\
Nov--Dec 2026 & Per-component fragility profile (C.1), two dataset pairs &
  Fragility table; figure set \\
Jan--Feb 2027 & Supervised oracle; repairs 1--2 with repeated seeds &
  Fraction-of-gap-closed per repair \\
Mar 2027 & Repairs 3--4; full FP16/INT8 matrix; Jetson latency for deployable repairs &
  Complete results \\
Apr 2027 & Writing; SJSU Research Showcase & Draft manuscript \\
May 2027 & Final report & Submitted by 30 May 2027 \\
\bottomrule
\end{tabular}
\end{center}

\textbf{Risk and contingency.} The oracle is the only step needing substantial GPU time;
without it, results are reported as absolute recovery rather than fraction-of-oracle ---
weaker but publishable. If C.2 refutes the road-geometry hypothesis, C.1 and C.3 are
unaffected.

% =================================================================== E
\section{E.\quad Anticipated Outcome}

\begin{enumerate}
\item \textbf{A per-component fragility profile} for roadside 3D detection across two dataset
      pairs --- attributing a cross-dataset drop to specific components rather than a single
      AP number.
\item \textbf{A measurement of how much of the gap each label-free repair recovers}, with
      confidence intervals, comparable to the fraction-of-gap convention in the adaptation
      literature.
\item \textbf{The first cross-dataset heading-error measurement for roadside sensors},
      reported as a metric distinct from AP.
\item \textbf{An answer on quantization}: whether INT8 amplifies domain fragility, measured on
      a mechanism-level metric sensitive enough to resolve it.
\item \textbf{Released artifacts} --- the TUMTraf\too DAIR converter, the JetPack~7 deployment
      recipe and all evaluation code, already public at
      \href{https://github.com/henrytran412/God-eye}{\texttt{github.com/henrytran412/God-eye}}.
\end{enumerate}

Target venues: SJSU Research Showcase; a workshop paper at CVPR, ICRA or IROS.

\textbf{Budget (\$300).} Portable SSD for dataset storage ($\sim$\$120); Jetson
power-monitoring hardware and cabling ($\sim$\$80); remainder for dataset access and
incidentals.

% ========================================================== references
\newpage
{\footnotesize
\begin{thebibliography}{99}
\setlength{\itemsep}{1.5pt}

\bibitem{ms3d} Tsai, D., Berrio, J.S., Shan, M., Nebot, E., Worrall, S.
  \emph{MS3D++: Ensemble of Experts for Multi-Source Unsupervised Domain Adaptation in 3D
  Object Detection.} IEEE T-IV, 2024. arXiv:2308.05988.

\bibitem{nuscenes} nuScenes detection benchmark --- mAOE, mATE, mASE definitions.
  nuScenes devkit. \url{https://github.com/nutonomy/nuscenes-devkit}

\bibitem{st3d} Yang, J., Shi, S., Wang, Z., Li, H., Qi, X.
  \emph{ST3D: Self-training for Unsupervised Domain Adaptation on 3D Object Detection.}
  CVPR 2021. arXiv:2103.05346.

\bibitem{st3dpp} Yang, J., Shi, S., Wang, Z., Li, H., Qi, X.
  \emph{ST3D++: Denoised Self-training for Unsupervised Domain Adaptation on 3D Object
  Detection.} IEEE T-PAMI 2022. arXiv:2108.06682.

\bibitem{cui} Cui, H., Chou, F.-C., Charland, J., Vallespi-Gonzalez, C., Djuric, N.
  \emph{Uncertainty-Aware Vehicle Orientation Estimation for Joint Detection-Prediction
  Models.} 2020. arXiv:2011.03114.

\bibitem{robust} Chandorkar, A.\ et al.
  \emph{Comprehensive Robustness Analysis of LiDAR-based 3D Object Detection in Autonomous
  Driving.} 2026. arXiv:2607.02074.

\bibitem{st3drepro} Independent reproducibility analysis of ST3D self-training. 2024.
  arXiv:2408.12708.

\bibitem{monodle} Ma, X., Zhang, Y., Xu, D., Zhou, D., Yi, S., Li, H., Ouyang, W.
  \emph{Delving into Localization Errors for Monocular 3D Object Detection.} CVPR 2021.
  arXiv:2103.16237.

\bibitem{gblobs} Mali\'{c}, D., Fruhwirth-Reisinger, C., Schulter, S., Possegger, H.
  \emph{GBlobs: Explicit Local Structure via Gaussian Blobs for Improved Cross-Domain
  LiDAR-based 3D Object Detection.} CVPR 2025. arXiv:2503.08639.

\bibitem{dair} Yu, H.\ et al.
  \emph{DAIR-V2X: A Large-Scale Dataset for Vehicle-Infrastructure Cooperative 3D Object
  Detection.} CVPR 2022.

\bibitem{tumtraf} Zimmer, W.\ et al.
  \emph{TUMTraf Intersection Dataset: All You Need for Urban 3D Camera-LiDAR Roadside
  Perception.} IEEE ITSC 2023.

\bibitem{nvidia} NVIDIA. \emph{CUDA-BEVFusion and CUDA-V2XFusion}, Lidar\_AI\_Solution.
  \url{https://github.com/NVIDIA-AI-IOT/Lidar_AI_Solution}

\end{thebibliography}}

\end{document}
